\documentclass[11pt]{article}

\usepackage[margin=0.85in]{geometry}
\usepackage{amsmath,amssymb,bm}
\usepackage{graphicx}
\usepackage{booktabs}
\usepackage{float}
\usepackage{xcolor}
\usepackage[hidelinks]{hyperref}

\title{Force Estimation, Realtime Filtering, and Leak-Free Uncertainty Bands}
\author{HSMR 2024 Effectiveness Force-Estimation Pipeline}
\date{\today}

\newcommand{\tauext}{\tau_{\mathrm{ext}}}
\newcommand{\taumeas}{\tau_{\mathrm{meas}}}
\newcommand{\taulstm}{\tau_{\mathrm{LSTM}}}
\newcommand{\Fhat}{\hat{\bm F}}
\newcommand{\tauhat}{\hat{\bm \tau}_{\mathrm{ext}}}

\begin{document}
\maketitle

\section{Purpose}

The goal is to estimate external wrench from motor torques and the robot model,
then attach a confidence interval (CI) without using the force sensor to tune or
compute that interval. The force sensor is used only for validation plots and
RMSE/coverage checks.

\begin{quote}
\textbf{No knowledge leakage rule.} The measured force sensor signal is never
used inside the uncertainty model. It is only used after estimation to compare
against \(\Fhat\).
\end{quote}

\section{Signals and Main Force Equation}

The LSTM predicts free-space joint torque. The measured motor torque contains
free-space dynamics plus external-contact effects. The torque residual is

\begin{equation}
\bm y_k
= \taumeas,k - \taulstm,k .
\end{equation}

Ideally,

\begin{equation}
\bm y_k \approx \tauext,k .
\end{equation}

Given a filtered external torque estimate \(\hat{\bm \tau}_{\mathrm{ext},k}\), the Cartesian wrench is
computed using the Jacobian transpose relationship:

\begin{equation}
\hat{\bm \tau}_{\mathrm{ext},k} = J_k^\mathsf{T} \Fhat_k ,
\end{equation}

\begin{equation}
\Fhat_k = J_k^{-\mathsf{T}} \hat{\bm \tau}_{\mathrm{ext},k}.
\end{equation}

The first three force channels are then rotated into the plotting frame using
the same rotation used in the original MATLAB script:

\begin{equation}
\Fhat_{k,1:3}^{\mathrm{plot}}
= R \Fhat_{k,1:3}.
\end{equation}

\section{Realtime Estimator Variants}

All estimators operate in joint-torque space first. The force projection
\(\Fhat = J^{-\mathsf{T}}\hat{\bm \tau}_{\mathrm{ext}}\) is applied after filtering.

\subsection{Direct Residual}

The simplest estimator uses the raw torque residual:

\begin{equation}
\hat{\bm \tau}_{\mathrm{ext},k} = \bm y_k .
\end{equation}

This is responsive but noisy and tends to preserve free-space false positives.

\subsection{Paper-Style Residual Observer}

The observer form discussed in the paper is

\begin{equation}
\dot{\bm r}
= K_i(\hat{\tau}_{\mathrm{ext}} - \bm r).
\end{equation}

In discrete time, the implementation uses an exact first-order update:

\begin{equation}
\bm r_k
= \alpha_k \bm r_{k-1}
+ (1-\alpha_k)\bm y_k,
\end{equation}

\begin{equation}
\alpha_k = \exp(-K_i\Delta t_k).
\end{equation}

The force estimate uses

\begin{equation}
\hat{\bm \tau}_{\mathrm{ext},k} = \bm r_k .
\end{equation}

This is realtime and stable. Larger \(K_i\) gives faster tracking but less
smoothing. The current tuned value is

\begin{equation}
K_i = 15~\mathrm{s}^{-1}.
\end{equation}

\subsection{Bias-Gated Observer}

The bias-gated observer tries to avoid false contact detection by estimating a
slow free-space bias only when the residual looks close to zero-contact:

\begin{equation}
\bm b_k \approx \text{slow bias in } \bm y_k .
\end{equation}

The corrected residual is

\begin{equation}
\tilde{\bm y}_k
= \operatorname{deadband}(\bm y_k - \bm b_k).
\end{equation}

Then the observer is applied to \(\tilde{\bm y}_k\). In our tests, nonzero bias
adaptation often worsened the result, so the practical setting has been

\begin{equation}
\text{bias gain} = 0.
\end{equation}

\subsection{Two-Time-Scale Observer}

This method subtracts a slow free-space estimate from a fast observer:

\begin{equation}
\hat{\bm \tau}_{\mathrm{ext},k}
= \operatorname{deadband}(\bm r^{\mathrm{fast}}_k
- \bm r^{\mathrm{slow}}_k).
\end{equation}

The idea is sensible, but in the current data it did not clearly outperform the
simpler tuned observer or the smooth-deadband methods.

\subsection{Smooth Torque-Space Deadband}

The most useful false-positive suppression came from applying a smooth
per-joint deadband before filtering. For joint \(i\),

\begin{equation}
d_i
= m \, s_i \, \sigma_{\tau,i},
\end{equation}

where \(m\) is the global deadband multiplier, \(s_i\) is a joint-specific
multiplier, and \(\sigma_{\tau,i}\) is estimated from the available torque
residual signal only.

The smooth gate is

\begin{equation}
g_i(y)
= \frac{1}
{1 + \exp\left(-\frac{|y|-d_i}{w_i}\right)},
\end{equation}

\begin{equation}
\tilde{y}_{i,k}
= y_{i,k} g_i(y_{i,k}).
\end{equation}

For the dVRK-Si robot, joints 1--2 are noisier and joint 3 tends to
underestimate in free space, so the current joint multipliers are

\begin{equation}
\bm s = [1,\;1,\;0.25,\;1,\;1,\;1].
\end{equation}

\subsection{Kalman Filter in Torque Space}

The Kalman filter is implemented before the force projection. The state is the
external joint torque:

\begin{equation}
\bm x_k = \tauext,k .
\end{equation}

The process model is a random walk:

\begin{equation}
\bm x_k = \bm x_{k-1} + \bm w_k,
\qquad
\bm w_k \sim \mathcal{N}(0,Q).
\end{equation}

The measurement is the LSTM torque residual:

\begin{equation}
\bm y_k
= \taumeas,k - \taulstm,k
= \bm x_k + \bm v_k,
\qquad
\bm v_k \sim \mathcal{N}(0,R).
\end{equation}

For each joint channel, the realtime Kalman recursion is

\begin{align}
P_{k|k-1} &= P_{k-1|k-1} + Q,\\
K_k &= \frac{P_{k|k-1}}{P_{k|k-1}+R},\\
\hat{x}_{k|k} &= \hat{x}_{k|k-1}
  + K_k(y_k-\hat{x}_{k|k-1}),\\
P_{k|k} &= (1-K_k)P_{k|k-1}.
\end{align}

The useful variant is not the raw Kalman filter alone, but
\[
\texttt{kalman\_smooth\_deadband},
\]
which first applies the smooth torque-space deadband and then applies the
Kalman filter. This keeps the false-positive suppression and replaces the fixed
observer gain with a probabilistic filter.

\section{Uncertainty Model}

The CI is leak-free: it uses motor torque residuals, the Jacobian, the estimated
force, and the filter disagreement. It does not use measured force.

\subsection{Torque-Side Noise Estimate}

Torque residual noise is estimated from first differences:

\begin{equation}
\Delta \bm y_k = \bm y_k - \bm y_{k-1}.
\end{equation}

The per-joint robust standard deviation is

\begin{equation}
\sigma_{\tau,i}
\approx
\frac{1.4826\,\operatorname{median}
\left(|\Delta y_i-\operatorname{median}(\Delta y_i)|\right)}
{\sqrt{2}}.
\end{equation}

The factor \(\sqrt{2}\) appears because differencing two independent noisy
samples doubles the variance.

\subsection{Propagating Torque Uncertainty to Force}

At each time step,

\begin{equation}
A_k = R_6 J_k^{-\mathsf{T}},
\end{equation}

where \(R_6\) applies the force-axis rotation to the first three wrench
components. The propagated wrench variance is

\begin{equation}
\bm \sigma_{\mathrm{prop},k}^2
= A_k^{\circ 2} \bm \sigma_\tau^{\,2},
\end{equation}

where \(A_k^{\circ 2}\) means elementwise squaring.

\subsection{Low-Force Ambiguity}

Near zero contact, the distinction between no contact, friction disturbance,
and small external force is ambiguous. This is modeled with

\begin{equation}
s_F(|\Fhat_k|)
= 1 + g_F \exp\left(-\frac{|\Fhat_k|}{F_0}\right).
\end{equation}

Current values:

\begin{equation}
g_F = 3.0,
\qquad
F_0 = 10~\mathrm{N}.
\end{equation}

\subsection{Observer/Direct Disagreement}

The direct projected wrench is

\begin{equation}
\Fhat^{\mathrm{direct}}_k
= J_k^{-\mathsf{T}}\bm y_k.
\end{equation}

The filtered estimate is

\begin{equation}
\Fhat^{\mathrm{filtered}}_k
= J_k^{-\mathsf{T}}\hat{\bm \tau}_{\mathrm{ext},k}.
\end{equation}

Their disagreement is used as a model-risk term:

\begin{equation}
\bm \sigma_{\mathrm{disagree},k}
= g_D
\left|
\Fhat^{\mathrm{direct}}_k
- \Fhat^{\mathrm{filtered}}_k
\right|.
\end{equation}

Current value:

\begin{equation}
g_D = 0.8.
\end{equation}

\subsection{Relative Force Uncertainty}

Large forces should not become unrealistically certain. A relative calibration
term is added:

\begin{equation}
\bm h_{\mathrm{rel},k}
= g_R |\Fhat_{k,1:3}|.
\end{equation}

Current value:

\begin{equation}
g_R = 0.10.
\end{equation}

The \(F_z\) channel also receives an empirical dVRK-Si axis correction:

\begin{equation}
h_{z,k} \leftarrow g_z h_{z,k},
\qquad
g_z = 2.0.
\end{equation}

\subsection{Final 95\% Band}

The final standard deviation estimate is

\begin{equation}
\bm \sigma_{F,k}
=
\sqrt{
\left(s_F(|\Fhat_k|)\bm\sigma_{\mathrm{prop},k}\right)^2
+
\bm\sigma_{\mathrm{disagree},k}^{\,2}
}.
\end{equation}

The 95\% half-width is

\begin{equation}
\bm h_k
= 1.96\,\bm\sigma_{F,k}.
\end{equation}

For the first three force channels, the relative term is combined as

\begin{equation}
\bm h_{k,1:3}
\leftarrow
\sqrt{
\bm h_{k,1:3}^{\,2}
+
\left(g_R|\Fhat_{k,1:3}|\right)^2
}.
\end{equation}

Then the plotted confidence interval is

\begin{equation}
\Fhat_k - \bm h_k
\leq
\bm F_k
\leq
\Fhat_k + \bm h_k.
\end{equation}

A maximum plotted half-width is used for readability:

\begin{equation}
h_{i,k} \leq 12~\mathrm{N}.
\end{equation}

\section{Current Key Hyperparameters}

\begin{table}[H]
\centering
\begin{tabular}{lll}
\toprule
Parameter & Current value & Meaning\\
\midrule
\(K_i\) & \(15~\mathrm{s}^{-1}\) & Observer speed\\
Deadband multiplier & \(20\) & Torque residual suppression\\
Joint deadband multipliers & \([1,1,0.25,1,1,1]\) & dVRK-Si joint-specific correction\\
Kalman process scale & \(0.08\) & Higher tracks faster but is noisier\\
Kalman measurement scale & \(1.0\) & Higher trusts residual less\\
\(g_F\) & \(3.0\) & Low-force uncertainty inflation\\
\(F_0\) & \(10~\mathrm{N}\) & Low-force decay scale\\
\(g_D\) & \(0.8\) & Observer/direct disagreement uncertainty\\
\(g_R\) & \(0.10\) & Relative force-scale uncertainty\\
\(g_z\) & \(2.0\) & \(F_z\) axis uncertainty correction\\
\bottomrule
\end{tabular}
\end{table}

\section{Validation Summary}

On the current dataset, the useful Kalman variant is
\[
\texttt{kalman\_smooth\_deadband}.
\]
Pure Kalman filtering of the raw residual is worse, because the raw residual
still contains free-space false positives and low-velocity friction/model
disturbances. The smooth deadband remains important.

\begin{table}[H]
\centering
\begin{tabular}{lrrrr}
\toprule
Method & \(F_x\) RMSE & \(F_y\) RMSE & \(F_z\) RMSE & Mean\\
\midrule
Raw Kalman, \(q=0.08\) & 2.4528 & 2.3832 & 1.4578 & 2.0979\\
Smooth-deadband observer & 1.6700 & 1.6260 & 1.4385 & 1.5782\\
Smooth-deadband Kalman, \(q=0.08\) & 1.6800 & 1.6266 & 1.4257 & 1.5775\\
\bottomrule
\end{tabular}
\end{table}

Interpretation: the Kalman-smooth method is essentially tied with the tuned
observer overall, slightly improves \(F_z\), and slightly worsens \(F_x\). Its
value is that the filter is more principled and tunable through \(Q/R\), while
remaining realtime.

\section{Plot Explanations}

\subsection{Force Estimate With 95\% Confidence Bands}

\begin{figure}[H]
\centering
\includegraphics[width=\textwidth]{outputs_velocity_context/external_force_ci.png}
\caption{Estimated force, force-sensor validation, and leak-free 95\% confidence
bands for the active estimator. The force sensor validates performance but is
not used to compute the band.}
\end{figure}

This plot answers whether the estimated force trajectory is plausible and
whether the CI captures most validation force samples. Not every point must be
covered by a 95\% interval. If whole contact events are outside the band, the
centerline estimator or uncertainty model needs improvement.

\subsection{CI Width Versus Estimated Contact Force}

\begin{figure}[H]
\centering
\includegraphics[width=0.85\textwidth]{outputs_velocity_context/ci_width_vs_contact_force.png}
\caption{Binned median CI half-width versus estimated force magnitude. This is
a diagnostic plot, not a forced monotonic law.}
\end{figure}

The original expectation was that uncertainty should shrink as force grows. The
current model is more nuanced:

\begin{itemize}
\item Near zero force, uncertainty is high because contact/no-contact is
ambiguous.
\item At moderate force, uncertainty can narrow.
\item At large force, uncertainty may widen because relative calibration/model
error grows with force magnitude.
\end{itemize}

Therefore, a mild U-shape or increasing high-force tail is expected when the
relative uncertainty term is active.

\subsection{Velocity Context for Force Uncertainty}

\begin{figure}[H]
\centering
\includegraphics[width=\textwidth]{outputs_velocity_context/velocity_with_force_uncertainty.png}
\caption{Tool speed from \(J\dot q\) shown with smoothed force CI half-widths.
This is a velocity-context plot, not a velocity confidence interval.}
\end{figure}

This plot helps address the advisor's concern: whether the force uncertainty is
only a proxy for velocity. The conclusion from the current data is that velocity
does affect some transient regions, especially fast teleoperation, but force CI
width is not explained by velocity alone.

\subsection{CI Width Versus Tool Speed}

\begin{figure}[H]
\centering
\includegraphics[width=0.85\textwidth]{outputs_velocity_context/ci_width_vs_velocity.png}
\caption{Binned median force CI half-width versus tool speed from \(J\dot q\).}
\end{figure}

If force uncertainty were purely velocity-driven, this plot would show a strong
monotonic increase. Instead, the relationship is weak and channel-dependent.
This suggests that force uncertainty comes from several sources: torque residual
ambiguity, filtering disagreement, Jacobian projection, and axis-dependent
model/calibration error.

\section{Main Conclusions}

\begin{enumerate}
\item The force-estimation pipeline is realtime: torque residual filtering is
causal, and force projection only requires the current Jacobian.
\item The Kalman filter is feasible and principled, but raw Kalman filtering is
not enough because the residual contains false-positive/free-space effects.
\item The best Kalman-style method is smooth torque deadband followed by the
Kalman filter.
\item The uncertainty band is leak-free because it never uses force-sensor
values during CI computation.
\item The CI is not simply a velocity uncertainty band. Velocity provides useful
context, but the force uncertainty is also driven by model/projection and
filter-disagreement terms.
\item A true velocity confidence interval would require a separate velocity
noise model or repeated velocity measurements. The current velocity plots are
diagnostics to disentangle velocity from force-estimation uncertainty.
\end{enumerate}

\end{document}
